\documentclass[12pt,oneside,a4paper]{article}

\usepackage[utf8]{inputenc}

\usepackage{amsfonts}
\usepackage{amsthm}        % Theorems
\usepackage{amsmath}

\newcommand{\bas}{\begin{eqnarray*}}
\newcommand{\eas}{\end{eqnarray*}}
\newcommand{\be}{\begin{equation}}
\newcommand{\ee}{\end{equation}}
\newcommand{\bea}{\begin{eqnarray}}
\newcommand{\eea}{\end{eqnarray}}

\title{Cassette Tape Counter and Playing Time}
\date{September 2026}
\author{Michael Jørgensen}

\begin{document}

\maketitle

\section{Problem}

On a cassette tape, the tape unwinds from one reel and winds onto the
other while passing the read head at a constant linear speed. Because
the amount of tape on each reel changes with time, the radius of each
wound pack --- and hence the angular velocity of each reel --- varies
during playback. Most tape decks have a mechanical or electronic
counter that records the total number of rotations of one of the
reels. We derive the relationship between the counter reading and the
elapsed playing time, and then discuss how to fit the resulting model
to experimental measurements.

\section{Physical model}

Let the tape move past the head with constant linear speed $v$. Track
the take-up reel, onto which tape winds starting from an (almost)
bare hub; the same argument applies to the supply reel with tape
unwinding, run in reverse.

Let
\begin{itemize}
\item $r_0$ be the radius of the bare hub,
\item $h$ be the effective thickness of a single wrap of tape (tape
  thickness together with whatever additional gap the winding
  process leaves between layers),
\item $R(t)$ be the outer radius of the wound tape pack at time $t$,
\item $\omega(t)$ be the angular velocity of the reel,
\item $N(t)$ be the counter reading, i.e.\ the total number of
  rotations completed by time $t$.
\end{itemize}

We assume the tape is thin, inextensible, and wound tightly with no
slip, so that the cross-sectional area of the pack grows in direct
proportion to the length of tape wound onto it.

\subsection{Radius as a function of time}

The length of tape wound onto the reel by time $t$ is $L(t) = vt$.
Modelling the pack as an annulus of tape of thickness $h$,
\be
\pi\bigl(R(t)^2 - r_0^2\bigr) = h\,L(t) = h v t,
\ee
so that
\be
R(t) = \sqrt{r_0^2 + \frac{hv}{\pi}\,t}\,.
\label{eq:radius}
\ee

\subsection{Angular velocity and total rotation}

The no-slip condition at the tape/pack interface means the rim speed
of the pack equals the tape speed, so
\be
\omega(t) = \frac{v}{R(t)} = \frac{v}{\sqrt{r_0^2 + \dfrac{hv}{\pi}t}}\,.
\ee
Integrating gives the total angle turned,
\bea
\theta(t) &=& \int_0^t \omega(t')\,dt' \nonumber\\
          &=& v\int_0^t \frac{dt'}{\sqrt{r_0^2 + \frac{hv}{\pi}t'}} \nonumber\\
          &=& \frac{2\pi}{h}\left(\sqrt{r_0^2+\frac{hv}{\pi}t} - r_0\right),
\eea
and the counter reading is $N(t) = \theta(t)/2\pi$:
\be
N(t) = \frac{1}{h}\left(\sqrt{r_0^2 + \frac{hv}{\pi}\,t}\; -\; r_0\right).
\label{eq:Noft}
\ee

\subsection{Time as a function of the counter reading}

Solving equation~(\ref{eq:Noft}) for $t$,
\be
\bigl(hN + r_0\bigr)^2 = r_0^2 + \frac{hv}{\pi}\,t,
\ee
which gives the key result:
\be
t(N) = \frac{2\pi r_0}{v}\,N \;+\; \frac{\pi h}{v}\,N^2.
\label{eq:tofN}
\ee
Playing time is a \emph{quadratic}, not linear, function of the
counter reading.

\subsection{Consistency checks}

Differentiating equation~(\ref{eq:tofN}),
\be
\frac{dt}{dN} = \frac{2\pi(r_0+hN)}{v} = \frac{2\pi R(t)}{v} = \frac{2\pi}{\omega(t)},
\ee
which is exactly the instantaneous rotation period at time $t$, as it
must be. At $N=0$, $dt/dN = 2\pi r_0/v$, the rotation period of the
bare hub --- correct, since the pack radius is $r_0$ at the start.
For large $N$ the quadratic term dominates and the counter advances
more slowly per unit time, matching the familiar observation that a
tape counter climbs quickly at first and more sluggishly later.

\section{Fitting the model to experimental data}

Equation~(\ref{eq:tofN}) is the model to be fitted to measured pairs
$(N_i, t_i)$.

\subsection{Choice of regression direction}

Two data-collection strategies are possible, and the correct
regression direction depends on which variable is actually measured
with negligible error and which one carries the noise.
\begin{itemize}
\item \textbf{Reading the counter at chosen times} (e.g.\ every 30 s):
  here $t$ is controlled by a clock and $N$ carries the noise
  (parallax, gear backlash). One should regress $N$ as a function of
  $t$ using the inverse form~(\ref{eq:Noft}), which is nonlinear in
  $r_0,h$ and requires an iterative (e.g.\ Levenberg--Marquardt)
  solver.
\item \textbf{Stopping the tape at chosen counter values} (e.g.\
  every 100 counts, timing with a stopwatch): here $N$ is controlled
  precisely and $t$ carries the noise. Equation~(\ref{eq:tofN}) is
  the correct regression direction, and, importantly, it is
  \emph{linear in the unknown parameters} $a = 2\pi r_0/v$ and
  $b = \pi h/v$, even though it is quadratic in $N$.
\end{itemize}
The second case is analytically simpler and is treated in detail
below; the first uses the same reasoning but requires nonlinear
least squares.

\subsection{Linear least squares for $t = aN + bN^2$}

Write
\be
t = c + aN + bN^2,
\label{eq:fitmodel}
\ee
including a possible constant offset $c$ to account for a
non-magnetic leader at the start of the tape or a counter not
exactly zeroed at the start of the recorded region. Physically
$c$ should be small and consistent with the leader length; if it is
statistically indistinguishable from what the leader implies (or
from zero, if there is no leader and the counter is zeroed exactly),
it may be dropped. Forcing $c=0$ without checking can bias $\hat a$
and $\hat b$.

Although equation~(\ref{eq:fitmodel}) is quadratic in $N$, it is
linear in the parameters $(c,a,b)$, since $N$ and $N^2$ can be
treated as two ordinary predictor variables. With data
$(N_i,t_i)$, $i=1,\dots,n$, form
\be
X=\begin{pmatrix}1 & N_1 & N_1^2\\ 1& N_2 & N_2^2\\ \vdots & \vdots& \vdots\\ 1& N_n & N_n^2\end{pmatrix},
\qquad
\mathbf t=\begin{pmatrix}t_1\\t_2\\\vdots\\t_n\end{pmatrix},
\ee
and solve the normal equations
\be
\hat{\boldsymbol\beta} = \begin{pmatrix}\hat c\\ \hat a\\ \hat b\end{pmatrix}
  = \bigl(X^{\mathsf T}X\bigr)^{-1}X^{\mathsf T}\mathbf t.
\ee
This is an ordinary polynomial regression, available directly in most
numerical packages (e.g.\ \texttt{numpy.linalg.lstsq},
\texttt{polyfit}, or a spreadsheet's \texttt{LINEST}).

\subsection{Uncertainty in the fitted coefficients}

Assume the timing errors $\varepsilon_i = t_i - \hat t_i$ are
independent and approximately Gaussian with common variance
$\sigma^2$ (reasonable if every $t_i$ is a stopwatch reading with
similar reaction-time error). Then
\be
\hat\sigma^2 = \frac{\mathrm{RSS}}{n-p}, \qquad
\mathrm{RSS} = \sum_i \bigl(t_i - \hat t_i\bigr)^2,
\ee
with $p$ the number of fitted parameters (2 or 3), and
\be
\mathrm{Cov}(\hat{\boldsymbol\beta}) = \hat\sigma^2 \bigl(X^{\mathsf T}X\bigr)^{-1}.
\ee
The diagonal entries give $\mathrm{Var}(\hat a)$ and
$\mathrm{Var}(\hat b)$; the off-diagonal entry gives
$\mathrm{Cov}(\hat a,\hat b)$, which is generally large in magnitude,
since $N$ and $N^2$ are strongly correlated over any single run of
data. Approximate $(1-\alpha)$ confidence intervals follow from the
Student-$t$ distribution with $n-p$ degrees of freedom,
$\hat a \pm t_{\alpha/2,\,n-p}\,\mathrm{SE}(\hat a)$, and similarly
for $\hat b$.

A narrow range of sampled $N$ values makes $N$ and $N^2$ nearly
proportional over that range, so $X^{\mathsf T}X$ becomes
ill-conditioned and $\hat a,\hat b$ individually become poorly
determined even though $\hat t(N)$ may still be well predicted
\emph{within} that range. Spreading measurements across the full
range of $N$ is the most effective way to control this.

\subsection{Recovering the physical parameters}

From the fitted coefficients,
\be
r_0 = \frac{av}{2\pi}, \qquad h = \frac{bv}{\pi}.
\ee
The timing data alone only ever constrain the two combinations
$a = 2\pi r_0/v$ and $b = \pi h/v$: there are three physical
unknowns $(r_0, h, v)$ but only two independent fitted numbers, so
the system is degenerate and one quantity must be fixed
independently. Two practical options are:
\begin{itemize}
\item take $v$ from the format specification (nominal compact
  cassette speed $v = 4.76\ \mathrm{cm/s} = 1\tfrac{7}{8}\
  \mathrm{in/s}$), or measure it directly by timing the full,
  known length of tape printed on the cassette shell --- this also
  checks whether the deck runs at nominal speed, since motor and
  capstan tolerances commonly cause a percent-level deviation;
\item measure $r_0$ directly with calipers on the bare hub (a
  purely geometric measurement with essentially no timing error),
  and/or take $h$ from the manufacturer's tape thickness
  specification, then solve for the remaining unknown(s), with the
  other equation serving as an internal consistency check.
\end{itemize}

If $v$ is taken as an exactly known constant, propagate only the fit
uncertainty:
\be
\sigma_{r_0} = \frac{v}{2\pi}\,\mathrm{SE}(\hat a), \qquad
\sigma_h = \frac{v}{\pi}\,\mathrm{SE}(\hat b).
\ee
If $v$ itself carries uncertainty $\sigma_v$ (e.g.\ from a motor
speed tolerance), combine the relative uncertainties in quadrature,
\be
\left(\frac{\sigma_{r_0}}{r_0}\right)^2
  = \left(\frac{\mathrm{SE}(\hat a)}{\hat a}\right)^2
  + \left(\frac{\sigma_v}{v}\right)^2,
\ee
and similarly for $h$ from $b$. When several independently measured
quantities are combined (e.g.\ a directly measured $r_0$ together
with the fitted $a,b$ to solve for $v$), a full delta-method
propagation using the Jacobian of $(r_0,h)$ with respect to
$(a,b,v)$, together with the $a$--$b$ covariance from
Section~3.3, is the rigorous approach.

\subsection{Sources of error and model validation}

\paragraph{Random errors} shrink with more or better data:
stopwatch reaction time (of order $0.1$--$0.3$ s per press, halved by
averaging repeated trials, or eliminated with an electronic timer);
counter quantization and backlash; and small capstan speed jitter
(wow and flutter).

\paragraph{Systematic errors} bias the fitted parameters and do not
shrink with more data of the same kind:
\begin{itemize}
\item \emph{Actual vs.\ nominal tape speed.} A deck running even a
  percent off nominal speed rescales $v$ and hence biases $r_0$ and
  $h$ together, though not $a$ and $b$ themselves. This is precisely
  the degeneracy of Section~3.4, and an independent speed check is
  the remedy.
\item \emph{Non-uniform pack density.} The model assumes a solid
  pack with uniform effective thickness $h$; uneven winding tension
  can make the effective thickness vary slowly with $N$. This shows
  up as systematic curvature in the residuals $t_i - \hat t_i$
  plotted against $N_i$: pure noise should look structureless, while
  a bowl or S-shape indicates the constant-$h$ assumption is
  breaking down, calling for a generalized model (e.g.\ $h$ varying
  with $N$, or an added cubic term).
\item \emph{Counter gearing ratio.} Many decks gear the counter off
  an idler rather than the reel itself, so the reading is
  $N = k \times (\text{reel rotations})$ for an unknown constant
  $k$. This preserves the quadratic form of $t(N)$ (rescaling $N$
  rescales $a$ and $b$ together) but means $r_0$ and $h$ are only
  recovered up to this ratio unless it is calibrated separately,
  e.g.\ by counting reel turns by hand against counter ticks over a
  short stretch of tape.
\item \emph{Leader tape and counter zero offset}, addressed via the
  intercept $c$ in Section~3.2.
\end{itemize}

\paragraph{Goodness of fit.} Report $R^2$ or, more usefully, the
residual standard deviation $\hat\sigma$ in the same units as the
timing precision, and compare it to the expected measurement noise;
a much larger $\hat\sigma$ signals a systematic effect rather than
random measurement error. Inspect the residual-vs-$N$ plot for
trends or curvature. If repeated measurements were taken at the same
$N$, their spread gives an independent estimate of pure measurement
noise, which can be compared to the regression residuals to
distinguish measurement noise from genuine model misspecification.

\subsection{Recommendations for experimental design}

\begin{itemize}
\item Spread the sampled $N$ values across the whole tape rather
  than clustering them near one end: points near $N=0$ anchor $a$
  (where $bN^2 \ll aN$), and points at large $N$ pin down $b$ (where
  the quadratic term dominates). This directly mitigates the
  multicollinearity issue of Section~3.3.
\item Replicate each measurement several times, rewinding to the
  same counter value, and average; this reduces stopwatch noise
  roughly as the inverse square root of the number of replicates.
\item Use the cassette's rated total playing time and tape length as
  an independent consistency check: $\hat t(N_{\max})$ should match
  the rated time, and the tape speed back-computed from length and
  time should match the assumed $v$.
\item If separate, accurate values of $r_0$ and $h$ are wanted
  rather than just the timing curve, obtain at least one independent
  direct measurement (most easily $r_0$, via calipers on the bare
  hub) rather than relying solely on the nominal tape speed --- this
  generally yields substantially better accuracy than extracting all
  three physical quantities from timing data alone.
\end{itemize}

\end{document}
